\documentclass[11pt,a4paper]{article}

% --- Codificación e idioma ---
\usepackage[utf8]{inputenc}
\usepackage[T1]{fontenc}
\usepackage[spanish,es-tabla]{babel}
\usepackage{lmodern}

% --- Geometría y formato general ---
\usepackage[margin=2.5cm]{geometry}
\usepackage{microtype}
\usepackage{parskip}
\usepackage{enumitem}
\setlist{topsep=2pt,itemsep=2pt,parsep=0pt}

% --- Tablas ---
\usepackage{array}
\usepackage{booktabs}
\usepackage{longtable}

% --- Diagramas simples ---
\usepackage{tikz}
\usetikzlibrary{positioning,arrows.meta}

% --- Cabeceras y numeración ---
\usepackage{fancyhdr}
\pagestyle{fancy}
\fancyhf{}
\fancyhead[L]{\small Bases de Datos II}
\fancyhead[R]{\small Resumen Teórico --- DCL (Práctico N.\textsuperscript{o} 2)}
\fancyfoot[C]{\thepage}
\renewcommand{\headrulewidth}{0.4pt}

% --- Color y código ---
\usepackage{xcolor}
\definecolor{codebg}{RGB}{248,248,248}
\definecolor{codeframe}{RGB}{200,200,200}
\definecolor{codekw}{RGB}{0,0,180}
\definecolor{codestr}{RGB}{163,21,21}
\definecolor{codecmt}{RGB}{0,128,0}
\definecolor{notebg}{RGB}{255,250,230}
\definecolor{noteframe}{RGB}{220,180,80}

\usepackage{listings}
\lstdefinelanguage{SQLext}[]{SQL}{
    morekeywords={NUMERIC,VARCHAR,DECIMAL,DATE,TIMESTAMP,TIME,
        BIGINT,VARCHAR2,NUMBER,ENUM,SCHEMA,SEQUENCE,TRIGGER,
        BEFORE,EACH,ROW,NEXTVAL,DUAL,RAISE,EXCEPTION,LOOP,
        IDENTIFIED,QUOTA,UNLIMITED,USERS,TABLESPACE,
        DATABASE,DOMAIN,GRANT,REVOKE,ROLE,USER,LOGIN,NOLOGIN,
        OPTION,REFERENCES,RESTRICT,CASCADE,
        AUTO\_INCREMENT,IF,EXISTS,ALTER,SHOW,USE,
        SEARCH\_PATH,EXECUTE,IMMEDIATE,PROFILE,PASSWORD,EXPIRE,
        CONNECT\_TIME,RESOURCE\_LIMIT,SYSTEM,SYSDBA,
        MAX\_CONNECTIONS\_PER\_HOUR,MAX\_QUERIES\_PER\_HOUR,
        MAX\_UPDATES\_PER\_HOUR,MAX\_USER\_CONNECTIONS,
        SESSIONS\_PER\_USER,IDLE\_TIME,
        ADMIN,GROUP,SUPERUSER,NOSUPERUSER,
        CREATEDB,NOCREATEDB,CREATEROLE,NOCREATEROLE,
        INHERIT,NOINHERIT,REPLICATION,NOREPLICATION,
        WITH,FROM,TO,ON,AS,AND,OR,IN,IS,NOT,NULL,
        INSERT,UPDATE,DELETE,SELECT,SHOW,REVOKE,GRANT,
        IDENTIFIED,BY,ENCRYPTED,UNENCRYPTED,VALID,UNTIL,
        CONNECTION,LIMIT,DEFAULT,TEMPORARY},
    sensitive=false,
    morecomment=[l]{--},
    morestring=[b]'
}
\lstset{
    language=SQLext,
    basicstyle=\ttfamily\footnotesize,
    keywordstyle=\color{codekw}\bfseries,
    stringstyle=\color{codestr},
    commentstyle=\color{codecmt}\itshape,
    backgroundcolor=\color{codebg},
    frame=single,
    rulecolor=\color{codeframe},
    breaklines=true,
    breakatwhitespace=true,
    columns=fullflexible,
    keepspaces=true,
    showstringspaces=false,
    captionpos=b,
    xleftmargin=0.4em,
    xrightmargin=0.4em,
    aboveskip=8pt,
    belowskip=8pt,
    tabsize=2
}

% --- Cajas de nota ---
\usepackage[most]{tcolorbox}
\newtcolorbox{nota}{colback=notebg,colframe=noteframe,arc=2pt,
    left=6pt,right=6pt,top=4pt,bottom=4pt,boxrule=0.6pt}

% --- Hipervínculos ---
\usepackage[hidelinks,bookmarks=true]{hyperref}

\title{Resumen Teórico --- DCL\\[4pt]
       \large Bases de Datos II --- Práctico N.\textsuperscript{o} 2}
\author{Tomás Rodeghiero}
\date{2026}

\begin{document}

% =====================================================
% Carátula
% =====================================================
\begin{titlepage}
    \centering
    \vspace*{1.5cm}

    {\large \textbf{Universidad Nacional de Río Cuarto}}\\[6pt]
    {\normalsize Facultad de Ciencias Exactas, Físico-Químicas y Naturales}\\[4pt]
    {\normalsize Departamento de Computación}\\[4pt]
    {\normalsize Licenciatura en Ciencias de la Computación}\\[3cm]

    {\Large \textbf{Bases de Datos II}}\\[1.2cm]

    {\Large \textbf{Resumen Teórico}}\\[6pt]
    {\large DCL --- Lenguaje de Control de Datos\\[2pt]
    Práctico N.\textsuperscript{o} 2}\\[3cm]

    \begin{tabular}{rl}
        \textbf{Alumno:} & Tomás Rodeghiero \\[4pt]
        \textbf{Año:}    & 2026 \\
    \end{tabular}

    \vfill
    {\small Material elaborado a partir del teórico
    \emph{Teorico\_3\_DCL}, el enunciado del práctico y las
    resoluciones realizadas.}
\end{titlepage}

\tableofcontents
\newpage

% =====================================================
\section{Introducción}
% =====================================================

El \textbf{DCL} (\emph{Data Control Language}) es la parte del lenguaje SQL
dedicada a la \textbf{seguridad} y al \textbf{control de acceso} a los
recursos de la base de datos. Mientras que el DDL define las estructuras
(tablas, vistas, dominios) y el DML manipula los datos, el DCL responde a
una pregunta distinta: \emph{¿quién puede hacer qué sobre cada objeto?}

Una base de datos en producción rara vez se accede con un único usuario.
Se necesita distinguir, por ejemplo, entre un \emph{cajero} que sólo puede
consultar y modificar ciertas tablas, un \emph{administrativo} con
permisos más amplios y un \emph{DBA} que controla la base completa.
Para sostener ese tipo de organización SQL provee tres mecanismos
complementarios:

\begin{itemize}
    \item \textbf{Usuarios}: cuentas con las que se inicia sesión en el motor.
    \item \textbf{Privilegios}: permisos individuales sobre objetos
    (tablas, vistas, esquemas, procedimientos, etc.).
    \item \textbf{Roles} (o \emph{grupos}): conjuntos nominados de
    privilegios que se asignan a uno o varios usuarios para evitar
    repetir trabajo de configuración.
\end{itemize}

\begin{nota}
\textbf{Idea clave del Teórico 3.} El estándar SQL define los conceptos de
\emph{usuario}, \emph{rol} y \emph{privilegio}, y los comandos
\texttt{GRANT} / \texttt{REVOKE}. Sin embargo, \emph{la administración
real de usuarios la implementa cada motor a su manera}: MySQL, PostgreSQL
y Oracle XE no son compatibles entre sí en este punto. Reconocer qué es
estándar y qué es específico es la mitad del práctico.
\end{nota}

% =====================================================
\section{Seguridad: el problema y los mecanismos}
% =====================================================

SQL ofrece la siguiente caja de herramientas para proteger los datos:

\begin{enumerate}
    \item \textbf{Autenticación}: definición de usuarios y roles, con
    sus contraseñas y, en algunos motores, restricción del \emph{host}
    desde el que se aceptan las conexiones.
    \item \textbf{Autorización}: asignación de privilegios sobre objetos.
    Sin un privilegio explícito un usuario no puede leer ni escribir
    nada que no le pertenezca.
    \item \textbf{Control de recursos}: límites de uso (cantidad de
    consultas por hora, tiempo de inactividad, número máximo de
    conexiones simultáneas) que evitan abusos.
    \item \textbf{Trazabilidad implícita}: el grafo de concesiones
    permite, en cualquier momento, contestar quién dio qué privilegio
    a quién.
\end{enumerate}

\paragraph{Modelo conceptual.} El teórico cierra con un diagrama que
sintetiza las relaciones entre los conceptos:

\begin{center}
\begin{tikzpicture}[node distance=1.6cm and 2.2cm,
    box/.style={rectangle,draw,minimum width=2.4cm,minimum height=0.8cm,
                font=\small},
    arr/.style={-{Latex},thick}]
    \node[box] (usu) {Usuario};
    \node[box,below=of usu]              (priv) {Privilegio};
    \node[box,right=of usu]              (rol)  {Rol};
    \node[box,right=of rol]              (per)  {Perfil};
    \node[box,below=of per]              (rec)  {Recurso};
    \draw[arr] (usu) -- node[left,font=\scriptsize]{1\;\;\;1-N} (priv);
    \draw[arr] (usu) -- node[above,font=\scriptsize]{1\;\;\;0-N} (rol);
    \draw[arr] (rol) -- node[right,font=\scriptsize]{1\;\;\;1-N} (priv);
    \draw[arr] (usu) -- node[above,font=\scriptsize,sloped]{1\;\;\;0-1} (per);
    \draw[arr] (per) -- node[right,font=\scriptsize]{1\;\;\;1-N} (rec);
\end{tikzpicture}
\end{center}

Un \emph{usuario} se asocia con uno o más \emph{privilegios}
(directamente o a través de \emph{roles}) y, opcionalmente, con un
\emph{perfil} que regula su consumo de \emph{recursos}.

% =====================================================
\section{Asignación de privilegios: \texttt{GRANT}}
% =====================================================

% -----------------------------------------------------
\subsection{Sintaxis general}
% -----------------------------------------------------

El comando \texttt{GRANT} concede privilegios sobre objetos a usuarios o
roles:

\begin{lstlisting}
GRANT {ALL | listaPrivilegios}
   ON TABLE listaTablas
   TO listaUsuariosRoles
[WITH GRANT OPTION];
\end{lstlisting}

Significado de los componentes:

\begin{itemize}
    \item \texttt{ALL}: otorga \emph{todos} los privilegios soportados
    sobre el objeto. Es cómodo, pero peligroso.
    \item \texttt{listaPrivilegios}: una o más operaciones permitidas.
    Las clásicas son \texttt{SELECT}, \texttt{INSERT}, \texttt{UPDATE},
    \texttt{DELETE}, \texttt{REFERENCES} y \texttt{EXECUTE}.
    \item \texttt{listaTablas}: el objeto (o lista de objetos) sobre el
    que se aplican; en el estándar es uno solo, varios motores aceptan
    listas o comodines (\texttt{*.*}, \texttt{db.*}).
    \item \texttt{listaUsuariosRoles}: a quién se le da el permiso.
    \item \texttt{WITH GRANT OPTION}: el receptor podrá, a su vez,
    delegar ese mismo permiso a un tercero.
\end{itemize}

% -----------------------------------------------------
\subsection{Privilegios habituales}
% -----------------------------------------------------

\begin{center}
\begin{tabular}{@{}l p{9.5cm}@{}}
\toprule
\textbf{Privilegio} & \textbf{Permite\dots} \\
\midrule
\texttt{SELECT}       & Leer filas de una tabla o vista. \\
\texttt{INSERT}       & Agregar filas. \\
\texttt{UPDATE}       & Modificar filas existentes (puede acotarse a
                        ciertas columnas). \\
\texttt{DELETE}       & Eliminar filas. \\
\texttt{REFERENCES}   & Crear claves foráneas y \texttt{CHECK}s que
                        referencien la tabla/columna. \\
\texttt{EXECUTE}      & Ejecutar funciones o procedimientos almacenados. \\
\texttt{CREATE}, \texttt{DROP}, \texttt{ALTER} & Operar sobre la
                        estructura de bases, tablas o índices. \\
\texttt{INDEX}        & Crear y eliminar índices. \\
\texttt{CREATE VIEW}, \texttt{SHOW VIEW} & Trabajar con vistas. \\
\texttt{CREATE ROUTINE}, \texttt{ALTER ROUTINE} & Crear/modificar
                        procedimientos almacenados. \\
\texttt{GRANT OPTION} & Delegar privilegios a otros usuarios. \\
\bottomrule
\end{tabular}
\end{center}

\begin{nota}
\textbf{Granularidad por columna.} Para \texttt{SELECT}, \texttt{INSERT},
\texttt{UPDATE} y \texttt{REFERENCES} es posible especificar la lista
de columnas a las que aplica el permiso, escribiéndolas entre paréntesis
después del nombre del privilegio: \texttt{GRANT SELECT(apellido,
nombre) ON personas TO U1}. Es la forma natural de cumplir consignas
del práctico como \emph{``puede consultar los campos apellido y nombre
de cliente''}.
\end{nota}

% -----------------------------------------------------
\subsection{Ejemplos básicos del teórico}
% -----------------------------------------------------

\begin{lstlisting}
-- 1) SELECT sobre una tabla a varios usuarios.
GRANT SELECT ON personas TO U1, U2;

-- 2) Permitir a U1 crear FKs hacia personas(dni).
GRANT REFERENCES (dni) ON personas TO U1;

-- 3) U1 podra, a su vez, delegar el SELECT.
GRANT SELECT ON personas TO U1 WITH GRANT OPTION;

-- 4) SELECT solo sobre dos columnas (PostgreSQL >= 9).
GRANT SELECT (apellido, nombre) ON personas TO U1;
\end{lstlisting}

% =====================================================
\section{Eliminación de privilegios: \texttt{REVOKE}}
% =====================================================

\texttt{REVOKE} es la operación inversa: retira privilegios previamente
otorgados.

\begin{lstlisting}
REVOKE {ALL | listaPrivilegios}
    ON listaTablas
  FROM listaUsuarios [RESTRICT | CASCADE];
\end{lstlisting}

\paragraph{\texttt{RESTRICT} vs \texttt{CASCADE}.} Cuando un usuario
otorgó un permiso (vía \texttt{WITH GRANT OPTION}) a otros, al revocárselo
hay que decidir qué hacer con esos terceros:

\begin{itemize}
    \item \texttt{CASCADE}: el revoke se propaga. Si U1 dio el permiso
    a U3 y a U1 se le revoca, U3 también lo pierde.
    \item \texttt{RESTRICT}: si revocar implicaría romper permisos
    transferidos, la operación falla. Permite preguntar antes de
    propagar.
\end{itemize}

\begin{lstlisting}
REVOKE SELECT ON personas FROM U1 RESTRICT;
\end{lstlisting}

\begin{nota}
\textbf{Por defecto, \texttt{REVOKE} es en cascada} en el estándar.
PostgreSQL implementa correctamente la diferencia
\texttt{RESTRICT}/\texttt{CASCADE}; \textbf{MySQL no la respeta}, y
revoca de forma directa sin distinguir explícitamente.
\end{nota}

% =====================================================
\section{Propietarios y permisos sobre el esquema}
% =====================================================

\textbf{Propiedad.} El \emph{propietario} de un objeto es siempre el
usuario que lo creó. Sólo el propietario (o un usuario con privilegios
administrativos suficientes) puede otorgar permisos sobre el objeto a
terceros. De ahí que en el práctico, antes de hacer cualquier
\texttt{GRANT}, se conecte el DBA o el dueño del esquema.

\textbf{Operaciones de esquema.} Las operaciones que modifican la
\emph{estructura} de la base (\texttt{CREATE}, \texttt{DROP},
\texttt{ALTER}) sólo pueden ser realizadas por el propietario del
esquema o por usuarios con los privilegios correspondientes
(\texttt{CREATE}, \texttt{DROP}, \texttt{ALTER} sobre la base).

% =====================================================
\section{Grafo de concesión de permisos}
% =====================================================

Una herramienta conceptual muy práctica para razonar sobre el estado
de los permisos es el \textbf{grafo de concesión}: un grafo dirigido
\emph{por recurso y tipo de privilegio} en el que los nodos son los
usuarios y las aristas representan concesiones.

\begin{itemize}
    \item Existe un arco $U_i \to U_j$ si $U_i$ concedió a $U_j$
    el privilegio en cuestión.
    \item La raíz es siempre el \emph{DBA} o el \emph{dueño del objeto}.
    \item \textbf{Un usuario tiene el permiso si y sólo si existe un
    camino desde la raíz hasta su nodo.}
\end{itemize}

Esa caracterización en términos de caminos es la que justifica el
comportamiento de \texttt{REVOKE \dots\ CASCADE}: si al cortar un arco
algunos nodos quedan sin camino desde la raíz, automáticamente pierden
el privilegio.

\paragraph{Ejemplo paso a paso (del Teórico 3).}

\textbf{(1) Estado inicial.} El DBA es el único con \texttt{UPDATE}
sobre \texttt{personas}.

\begin{center}
\begin{tikzpicture}[node distance=1cm and 1.6cm,
    n/.style={circle,draw,minimum size=8mm,inner sep=0pt,font=\small},
    arr/.style={-{Latex}}]
    \node[n] (dba) {DBA};
    \node[n,right=of dba,yshift=0.6cm]  (u1) {U1};
    \node[n,right=of dba,yshift=-0.6cm] (u2) {U2};
    \node[n,right=2cm of u1]            (u3) {U3};
\end{tikzpicture}
\end{center}

\textbf{(2)} El DBA ejecuta:
\texttt{GRANT UPDATE ON personas TO U1, U2 WITH GRANT OPTION}.

\begin{center}
\begin{tikzpicture}[node distance=1cm and 1.6cm,
    n/.style={circle,draw,minimum size=8mm,inner sep=0pt,font=\small},
    arr/.style={-{Latex}}]
    \node[n] (dba) {DBA};
    \node[n,right=of dba,yshift=0.6cm]  (u1) {U1};
    \node[n,right=of dba,yshift=-0.6cm] (u2) {U2};
    \node[n,right=2cm of u1]            (u3) {U3};
    \draw[arr] (dba) -- (u1);
    \draw[arr] (dba) -- (u2);
\end{tikzpicture}
\end{center}

\textbf{(3)} U1 ejecuta:
\texttt{GRANT UPDATE ON personas TO U3}.

\begin{center}
\begin{tikzpicture}[node distance=1cm and 1.6cm,
    n/.style={circle,draw,minimum size=8mm,inner sep=0pt,font=\small},
    arr/.style={-{Latex}}]
    \node[n] (dba) {DBA};
    \node[n,right=of dba,yshift=0.6cm]  (u1) {U1};
    \node[n,right=of dba,yshift=-0.6cm] (u2) {U2};
    \node[n,right=2cm of u1]            (u3) {U3};
    \draw[arr] (dba) -- (u1);
    \draw[arr] (dba) -- (u2);
    \draw[arr] (u1)  -- (u3);
\end{tikzpicture}
\end{center}

\textbf{(4)} El DBA ejecuta:
\texttt{REVOKE UPDATE ON personas FROM U1 CASCADE}.
Como el único camino al nodo U3 pasaba por U1, U3 también pierde el
privilegio.

\begin{center}
\begin{tikzpicture}[node distance=1cm and 1.6cm,
    n/.style={circle,draw,minimum size=8mm,inner sep=0pt,font=\small},
    arr/.style={-{Latex}}]
    \node[n] (dba) {DBA};
    \node[n,right=of dba,yshift=0.6cm]  (u1) {U1};
    \node[n,right=of dba,yshift=-0.6cm] (u2) {U2};
    \node[n,right=2cm of u1]            (u3) {U3};
    \draw[arr] (dba) -- (u2);
\end{tikzpicture}
\end{center}

% =====================================================
\section{Roles}
% =====================================================

Un \textbf{rol} representa un \emph{tipo de usuario}: agrupa privilegios
asociados a una función dentro de la organización (cajero, asesor,
empleado, encargado). En lugar de repetir las mismas concesiones para
cada cuenta, se declara un rol una sola vez y se le asigna a las cuentas
que lo necesiten. Cuando aparece un nuevo cajero, basta con asignarle
el rol \emph{cajero} y queda configurado.

\paragraph{Administración de roles.}

\begin{lstlisting}
-- Crear y eliminar roles.
CREATE ROLE nombreRole;
DROP ROLE nombreRole;

-- Conceder un rol a un usuario.
GRANT role TO usuario;

-- Conceder un rol a otro rol (anidamiento).
GRANT role1 TO role2;
\end{lstlisting}

Las dos últimas formas explican por qué \texttt{GRANT} aparece dos veces
en una secuencia típica de configuración: primero se otorgan los
privilegios al rol y luego se otorga el rol a los usuarios.

% =====================================================
\section{Diferencias entre motores}
% =====================================================

La administración de usuarios y roles \emph{no es estándar}. Cada motor
tiene sus comandos, su modelo de cuentas y sus opciones de seguridad.

% -----------------------------------------------------
\subsection{MySQL}
% -----------------------------------------------------

MySQL trabaja con \emph{cuentas} de la forma \texttt{usuario@host}, donde
\emph{host} indica desde qué máquina se acepta la conexión:

\begin{itemize}
    \item \texttt{'app'@'localhost'}: sólo conexiones locales.
    \item \texttt{'app'@'\%'}: el comodín \texttt{\%} acepta cualquier
    host.
    \item Un mismo nombre de usuario puede existir con distintos hosts y
    distintas contraseñas.
\end{itemize}

\paragraph{Creación y borrado de cuentas.}

\begin{lstlisting}
CREATE USER 'usuario'@'host' [IDENTIFIED BY 'password'];
DROP USER 'usuario'@'host';
\end{lstlisting}

\paragraph{Roles (desde MySQL 8).}

\begin{lstlisting}
CREATE ROLE [IF NOT EXISTS] role [, role] ...;
GRANT role [, role] ... TO user_or_role [, user_or_role] ...
    [WITH ADMIN OPTION];
\end{lstlisting}

\paragraph{Límites de recursos.} Desde MySQL 5.0 se pueden acotar:

\begin{itemize}
    \item \texttt{MAX\_QUERIES\_PER\_HOUR}
    \item \texttt{MAX\_UPDATES\_PER\_HOUR}
    \item \texttt{MAX\_CONNECTIONS\_PER\_HOUR}
    \item \texttt{MAX\_USER\_CONNECTIONS} (conexiones simultáneas).
\end{itemize}

\begin{lstlisting}
-- Forma clasica (hasta 5.7) con GRANT.
GRANT ALL ON *.* TO 'app'@'localhost' WITH
    MAX_QUERIES_PER_HOUR 20
    MAX_UPDATES_PER_HOUR 10
    MAX_CONNECTIONS_PER_HOUR 5
    MAX_USER_CONNECTIONS 2;

-- Forma actual (>= 8) con ALTER USER.
ALTER USER 'app'@'localhost'
    WITH MAX_QUERIES_PER_HOUR 500
         MAX_UPDATES_PER_HOUR 100;
\end{lstlisting}

\begin{nota}
\textbf{Para MySQL 8+.} Los límites de recursos en \texttt{GRANT} fueron
desaprobados como mecanismo principal y la cátedra recomienda usar
\texttt{ALTER USER}. La sintaxis \texttt{GRANT \dots\ WITH MAX\_\dots}
sigue funcionando por compatibilidad.
\end{nota}

% -----------------------------------------------------
\subsection{PostgreSQL}
% -----------------------------------------------------

En PostgreSQL los conceptos de \emph{usuario}, \emph{rol} y
\emph{grupo} son lo mismo: todos son \emph{roles}. La diferencia es
operativa:

\begin{itemize}
    \item Un rol con permiso de \texttt{LOGIN} se comporta como
    \emph{usuario}.
    \item Un rol \texttt{NOLOGIN} se comporta como \emph{grupo}.
    Los usuarios se hacen miembros de él para heredar sus privilegios.
\end{itemize}

\paragraph{Creación de roles.} La sintaxis es muy expresiva:

\begin{lstlisting}
CREATE ROLE nombre [WITH] [option ...];
-- option puede ser, entre otras:
--   SUPERUSER | NOSUPERUSER
--   CREATEDB  | NOCREATEDB
--   CREATEROLE| NOCREATEROLE
--   INHERIT   | NOINHERIT
--   LOGIN     | NOLOGIN
--   CONNECTION LIMIT n
--   {ENCRYPTED | UNENCRYPTED} PASSWORD 'pwd'
--   VALID UNTIL 'timestamp'
--   IN ROLE/IN GROUP rolename, ...
--   ROLE rolename, ...
\end{lstlisting}

Existe además el atajo \texttt{CREATE USER} (que es \texttt{CREATE ROLE}
con \texttt{LOGIN} por defecto) y la utilidad de línea de comandos
\texttt{createuser}, con flags como \texttt{-P} (pide contraseña),
\texttt{-d} (puede crear bases), \texttt{-A} (no puede crear
usuarios) y \texttt{-h} (host).

\paragraph{Grupos.} \texttt{CREATE GROUP} es un alias histórico de
\texttt{CREATE ROLE} sin \texttt{LOGIN}:

\begin{lstlisting}
CREATE GROUP groupName [WITH SYSID groupid] [USER username, ...];
ALTER GROUP groupName {ADD | DROP} USER username;  -- equivale a GRANT/REVOKE
DROP GROUP groupName;
\end{lstlisting}

\paragraph{Acceso desde clientes.} Por defecto PostgreSQL acepta
conexiones sólo desde \texttt{localhost}. Para permitir clientes
remotos hay que editar \texttt{pg\_hba.conf}, por ejemplo:

\begin{lstlisting}
# TYPE  DATABASE  USER  CIDR-ADDRESS  METHOD
host    all       all   0.0.0.0/0     md5
\end{lstlisting}

% -----------------------------------------------------
\subsection{Oracle XE}
% -----------------------------------------------------

En Oracle XE \textbf{crear un usuario equivale a crear un esquema}: el
usuario es el dueño de su propio espacio de objetos.

\paragraph{Creación de usuario (ejemplo).}

\begin{lstlisting}
CREATE USER usuario
    IDENTIFIED BY usuario
    DEFAULT TABLESPACE users
    QUOTA 10M ON users
    TEMPORARY TABLESPACE temp
    QUOTA 5M ON system
    PASSWORD EXPIRE;
\end{lstlisting}

\texttt{PASSWORD EXPIRE} fuerza al usuario a cambiar la clave en su
próximo \emph{login}, recurso muy usado para cuentas recién creadas
(es justamente lo que pide el inciso 4 del práctico para
\emph{empleado1}).

\paragraph{Roles predefinidos.} Oracle XE trae tres roles listos para
asignar:

\begin{itemize}
    \item \texttt{CONNECT}: conectarse a la base de datos.
    \item \texttt{RESOURCE}: para desarrolladores; permite crear
    objetos pero no vistas.
    \item \texttt{DBA}: rol de administrador.
\end{itemize}

\paragraph{Perfiles (\texttt{PROFILE}).} Oracle permite definir
\emph{perfiles} que limitan el uso de recursos por usuario y se
asignan luego con \texttt{ALTER USER}:

\begin{lstlisting}
CREATE PROFILE desarrolador LIMIT
    SESSIONS_PER_USER 1
    IDLE_TIME 30
    CONNECT_TIME 600;
\end{lstlisting}

Es la forma natural de cumplir consignas como ``\emph{el director sólo
puede estar conectado un máximo de 20 minutos}'' del Ejercicio 4.

\paragraph{Cuentas APEX vs cuentas del motor.} Oracle XE distingue las
cuentas de la aplicación web APEX (creadas dentro de un \emph{workspace})
de los usuarios del motor de base de datos. \textbf{Los usuarios APEX
no son usuarios de la base}; sí lo son los esquemas que esos workspaces
crean. El DBA del motor es \texttt{SYS} o \texttt{SYSTEM}, y la
administración por shell se recomienda con \texttt{sqlplus}.

\paragraph{Comparación rápida entre motores.}

\begin{center}
\begin{tabular}{@{}p{3.4cm} p{3.5cm} p{3.5cm} p{3.6cm}@{}}
\toprule
\textbf{Tema} & \textbf{MySQL 8+} & \textbf{PostgreSQL} &
\textbf{Oracle XE} \\
\midrule
Cuenta & \texttt{usuario@host} con password & \emph{Role} con
\texttt{LOGIN} & Usuario equivale a esquema. \\
\addlinespace
Roles & \texttt{CREATE ROLE} (desde 8) & Todo es rol;
\texttt{CREATE ROLE} & \texttt{CREATE ROLE}; trae
\texttt{CONNECT}/\texttt{RESOURCE}/\texttt{DBA}. \\
\addlinespace
Restricción de host & Parte del nombre de cuenta &
\texttt{pg\_hba.conf} & Configuración de red separada. \\
\addlinespace
Limites de recursos & \texttt{MAX\_*\_PER\_HOUR} & \texttt{CONNECTION
LIMIT} en el rol & \texttt{CREATE PROFILE} +
\texttt{ALTER USER \dots\ PROFILE}. \\
\addlinespace
\texttt{REVOKE} con \texttt{RESTRICT} / \texttt{CASCADE} & No respeta
realmente la diferencia & Implementado del estándar & Comportamiento
estándar. \\
\addlinespace
Forzar cambio de clave & \texttt{ALTER USER \dots\ PASSWORD EXPIRE} &
\texttt{VALID UNTIL} & \texttt{PASSWORD EXPIRE} en \texttt{CREATE} o
\texttt{ALTER USER}. \\
\bottomrule
\end{tabular}
\end{center}

% =====================================================
\section{Administración: usuarios, contraseñas y herramientas}
% =====================================================

\paragraph{MySQL.} La instalación crea el superusuario \texttt{root};
desde la versión 5 solicita la contraseña durante la instalación, y
desde la 8 en Linux por defecto autentica con el sistema operativo. Las
tablas que almacenan usuarios y permisos viven en la base interna
\texttt{mysql}. Para la administración se recomiendan
\textbf{MySQL Workbench} (interfaz actual), \texttt{phpMyAdmin}
(servicios de hosting) o \texttt{MySQLAdministrator} (desactualizado).

\paragraph{PostgreSQL.} La instalación crea el superusuario
\texttt{postgres} (con un usuario homónimo a nivel de sistema
operativo en Linux). Para conectarse desde aplicaciones cliente con
autenticación por base de datos --- y no por SO --- hay que ajustar
\texttt{pg\_hba.conf}. Herramientas habituales:
\textbf{PGAdmin IV} y \texttt{phpPgAdmin}.

\paragraph{Oracle XE.} Se administra por web en
\texttt{http://localhost:8080/apex} o, recomendado, por consola con
\texttt{sqlplus}. En 11XE el usuario DBA del motor es \texttt{SYS} o
\texttt{SYSTEM} y el workspace por defecto es \texttt{internal}.

% =====================================================
\section{Conexión con los ejercicios del práctico}
% =====================================================

El Práctico 2 está diseñado para ejercitar exactamente los conceptos
anteriores. Cada ejercicio toma una de las bases de datos del Práctico 1
y agrega usuarios y permisos sobre ella.

% -----------------------------------------------------
\subsection*{Ejercicio 1 (MySQL) --- base de Cliente/Producto/Factura}
% -----------------------------------------------------

Sobre la base \texttt{practico1a} se crean los usuarios
\emph{vendedor1}, \emph{vendedor2} y \emph{administrador} con permisos
muy diferenciados. La consigna combina:

\begin{itemize}
    \item \texttt{INSERT} sobre una tabla específica (vendedor1 sobre
    \emph{cliente}; vendedor2 sobre \emph{factura}).
    \item \texttt{SELECT} acotado a una tabla (vendedor2 sobre
    \emph{producto}).
    \item \texttt{UPDATE} a nivel de \textbf{columna} con
    \texttt{WITH GRANT OPTION} (administrador sobre
    \emph{producto.descripcion}).
    \item \texttt{DELETE} colectivo sobre \emph{cliente} para los tres.
    \item \texttt{REVOKE ALL} para limpiar al final los permisos del
    vendedor2.
\end{itemize}

\paragraph{Esqueleto típico.}

\begin{lstlisting}
CREATE USER 'vendedor1'@'localhost' IDENTIFIED BY 'pwd';
CREATE USER 'vendedor2'@'localhost' IDENTIFIED BY 'pwd';
CREATE USER 'administrador'@'localhost' IDENTIFIED BY 'pwd';

USE practico1a;

GRANT INSERT ON practico1a.cliente   TO 'vendedor1'@'localhost';

GRANT INSERT ON practico1a.factura   TO 'vendedor2'@'localhost';
GRANT SELECT ON practico1a.producto  TO 'vendedor2'@'localhost';

GRANT UPDATE (descripcion) ON practico1a.producto
    TO 'administrador'@'localhost' WITH GRANT OPTION;

-- DELETE para los tres.
GRANT DELETE ON practico1a.cliente TO
    'vendedor1'@'localhost',
    'vendedor2'@'localhost',
    'administrador'@'localhost';

-- Verificacion y limpieza final.
SHOW GRANTS FOR 'vendedor2'@'localhost';
REVOKE ALL PRIVILEGES, GRANT OPTION FROM 'vendedor2'@'localhost';
\end{lstlisting}

% -----------------------------------------------------
\subsection*{Ejercicio 2 (MySQL) --- base de Vehículos/Estacionamiento}
% -----------------------------------------------------

Aquí entra en juego la dimensión \emph{host}: el
\emph{encargado\_estacionamiento} se conecta sólo desde
\texttt{localhost} y el \emph{cobrador} desde cualquier máquina
(\texttt{\%}). El cobrador, además, queda limitado en recursos
con \texttt{MAX\_CONNECTIONS\_PER\_HOUR=3} y
\texttt{MAX\_QUERIES\_PER\_HOUR=10}.

\begin{lstlisting}
CREATE USER 'encargado_estacionamiento'@'localhost' IDENTIFIED BY 'pwd';
CREATE USER 'cobrador'@'%' IDENTIFIED BY 'pwd';

GRANT SELECT, INSERT ON practico1b.parquimetro
   TO 'encargado_estacionamiento'@'localhost';
GRANT SELECT ON practico1b.*
   TO 'encargado_estacionamiento'@'localhost';

GRANT SELECT ON practico1b.estacionamiento TO 'cobrador'@'%';

ALTER USER 'cobrador'@'%' WITH
    MAX_CONNECTIONS_PER_HOUR 3
    MAX_QUERIES_PER_HOUR     10;

-- Inciso d).
REVOKE SELECT ON practico1b.estacionamiento FROM 'cobrador'@'%';
\end{lstlisting}

% -----------------------------------------------------
\subsection*{Ejercicio 3 (PostgreSQL) --- base de Accidentes}
% -----------------------------------------------------

Es el ejercicio que introduce \textbf{roles compartidos}. El rol
\emph{empleado} consolida lo común a todos (\texttt{SELECT} sobre
\texttt{cliente(apellido, nombre)}); a partir de ahí cada usuario
recibe los permisos específicos:

\begin{itemize}
    \item \emph{asesor}: \texttt{SELECT} en \emph{taller} e
    \texttt{INSERT} en \emph{accidente}.
    \item \emph{administrativo}: \texttt{SELECT} en \emph{automóvil}.
    \item \emph{encargado}: \texttt{DROP} sobre todas las tablas y
    \texttt{UPDATE} sobre \texttt{categoria(tasa)}.
\end{itemize}

\begin{lstlisting}
-- Rol comun.
CREATE ROLE empleado;
GRANT SELECT (apellido, nombre) ON cliente TO empleado;

-- Usuarios reales.
CREATE ROLE asesor         LOGIN PASSWORD 'pwd' IN ROLE empleado;
CREATE ROLE administrativo LOGIN PASSWORD 'pwd' IN ROLE empleado;
CREATE ROLE encargado      LOGIN PASSWORD 'pwd' IN ROLE empleado;

GRANT SELECT ON taller TO asesor;
GRANT INSERT ON accidente TO asesor;
GRANT SELECT ON automovil TO administrativo;
GRANT UPDATE (tasa) ON categoria TO encargado;
-- DROP no es un privilegio que se otorgue: se concede ownership o
-- se ejecuta como propietario / superuser.
\end{lstlisting}

% -----------------------------------------------------
\subsection*{Ejercicio 4 (Oracle) --- base de Accidentes}
% -----------------------------------------------------

El Ejercicio 4 cierra el práctico ejercitando los aspectos más
\emph{Oracle}: rol explícito, perfil para acotar tiempo de conexión,
\texttt{PASSWORD EXPIRE} y \texttt{WITH GRANT OPTION} a nivel de
columnas.

\begin{lstlisting}
-- Rol y privilegios comunes.
CREATE ROLE empleados;
GRANT SELECT ON automovil TO empleados;
GRANT SELECT ON categoria TO empleados;

-- Perfil para limitar al director.
CREATE PROFILE perfil_director LIMIT
    CONNECT_TIME 20;

-- Usuarios.
CREATE USER empleado1 IDENTIFIED BY pwd PASSWORD EXPIRE;
CREATE USER empleado2 IDENTIFIED BY pwd;
CREATE USER director  IDENTIFIED BY pwd PROFILE perfil_director;

GRANT CONNECT TO empleado1, empleado2, director;
GRANT empleados TO empleado1, empleado2;

-- Privilegios especificos.
GRANT INSERT, UPDATE, DELETE ON accidente TO empleado1;
GRANT SELECT (apellido, nombre, direccion) ON cliente
    TO empleado2 WITH GRANT OPTION;

-- Director: SELECT sobre todo.
GRANT SELECT ON cliente   TO director;
GRANT SELECT ON automovil TO director;
GRANT SELECT ON categoria TO director;
GRANT SELECT ON taller    TO director;
GRANT SELECT ON accidente TO director;
\end{lstlisting}

% =====================================================
\section{Verificación de privilegios}
% =====================================================

Siempre que se asignan permisos hay que verificarlos. Cada motor tiene
su forma:

\begin{itemize}
    \item \textbf{MySQL}: \texttt{SHOW GRANTS FOR 'usuario'@'host';}
    \item \textbf{PostgreSQL}: las vistas de catálogo
    \texttt{information\_schema.table\_privileges},
    \texttt{column\_privileges} y \texttt{role\_table\_grants}; en
    \texttt{psql}, los meta-comandos \texttt{\textbackslash dp} y
    \texttt{\textbackslash z}.
    \item \textbf{Oracle}: las vistas \texttt{USER\_TAB\_PRIVS},
    \texttt{USER\_ROLE\_PRIVS}, \texttt{USER\_SYS\_PRIVS} y, para los
    DBA, sus equivalentes \texttt{DBA\_*}.
\end{itemize}

Una buena práctica del práctico es \emph{loguearse efectivamente con
cada usuario creado} y probar las operaciones permitidas (deben
funcionar) y las prohibidas (deben fallar con error de privilegio).
Si los \emph{tests} positivos y negativos coinciden con la consigna,
los permisos están bien definidos.

% =====================================================
\section{Ideas clave para repasar antes del parcial}
% =====================================================

\begin{enumerate}
    \item \textbf{DCL es el sublenguaje de seguridad.} Sus dos comandos
    centrales son \texttt{GRANT} y \texttt{REVOKE}.
    \item \textbf{Estándar vs implementación.} El estándar SQL define
    usuarios, roles y privilegios; la administración real la implementa
    cada motor. Hay que saber qué partes son comunes y qué partes
    cambian.
    \item \textbf{Privilegios típicos}: \texttt{SELECT}, \texttt{INSERT},
    \texttt{UPDATE}, \texttt{DELETE}, \texttt{REFERENCES},
    \texttt{EXECUTE}, \texttt{CREATE}, \texttt{DROP}, \texttt{ALTER},
    \texttt{INDEX}, \texttt{CREATE VIEW}.
    \item \textbf{Granularidad}: los privilegios pueden ser globales
    (\texttt{*.*}), por base, por tabla o por columna. Para columnas se
    usa la lista entre paréntesis después del nombre del privilegio.
    \item \textbf{\texttt{WITH GRANT OPTION}} habilita la
    \emph{delegación}: el receptor puede a su vez otorgar el privilegio
    a otros.
    \item \textbf{\texttt{REVOKE \dots\ CASCADE / RESTRICT}} controla
    qué pasa con los privilegios delegados al revocar uno de ellos.
    Por defecto el comportamiento es en cascada.
    \item \textbf{Roles} son ``tipos de usuario''. Se crean con
    \texttt{CREATE ROLE}, se les otorgan privilegios con \texttt{GRANT
    \dots\ ON \dots\ TO role} y se asignan a usuarios con
    \texttt{GRANT role TO usuario}.
    \item \textbf{El propietario} de un objeto es siempre quien lo
    creó. Sólo él (o un DBA) puede otorgar privilegios sobre él.
    \item \textbf{Grafo de concesión}: existe un permiso si existe un
    camino desde la raíz (DBA o dueño) al usuario. Si al revocar se
    rompe el último camino, el usuario lo pierde.
    \item \textbf{MySQL} maneja cuentas como \texttt{usuario@host} y
    permite limitar recursos con \texttt{MAX\_QUERIES\_PER\_HOUR},
    \texttt{MAX\_CONNECTIONS\_PER\_HOUR}, etc.
    \item \textbf{PostgreSQL} unifica usuarios, roles y grupos: todos
    son \emph{roles}. \texttt{LOGIN} indica que pueden iniciar sesión.
    \item \textbf{Oracle} asocia usuario y esquema; ofrece
    \emph{perfiles} (\texttt{CREATE PROFILE}) para limitar recursos
    y \texttt{PASSWORD EXPIRE} para forzar cambio de clave.
    \item \textbf{Verificar siempre}: \texttt{SHOW GRANTS} en MySQL,
    \texttt{\textbackslash dp} en PostgreSQL, \texttt{USER\_*\_PRIVS}
    en Oracle. Y, además, probar conectándose como cada usuario.
\end{enumerate}

% =====================================================
\section*{Resumen final}
% =====================================================
\addcontentsline{toc}{section}{Resumen final}

DCL es la pieza que convierte una base de datos correctamente
estructurada (DDL) y poblada (DML) en un sistema realmente
\emph{multiusuario} y seguro. La idea central es simple: cada
operación debe estar respaldada por un privilegio explícito sobre el
objeto correspondiente, y esos privilegios pueden viajar a través de
roles (para escalar a muchos usuarios), de \texttt{WITH GRANT OPTION}
(para delegar) y de un grafo de concesiones que ordena las relaciones
entre quienes otorgan y reciben.

Si quedan claros los cinco bloques --- usuarios, privilegios, roles,
\texttt{WITH GRANT OPTION} con \texttt{REVOKE \dots\ CASCADE/RESTRICT}
y la administración por motor --- el Práctico 2 se vuelve un mapeo
directo de la teoría: cada inciso se traduce en una secuencia
\emph{crear usuario / otorgar privilegio / verificar / revocar}. El
verdadero examen no es escribir el comando, sino justificar por qué
ese conjunto de permisos cumple la consigna ni más ni menos que lo
necesario.

\end{document}
